\subsection{平坦模的例子}

\begin{enumerate}
\def\labelenumi{(\arabic{enumi})}
\item
  对任意环 A， \(A_{d}\) 显然是平坦 A-模（《代数学》第二章 §3 第 4 目命题 4）。于是由第 3 目命题 2 (i) 得：每个自由右 A-模，以及更一般地每个投射右 A-模（《代数学》第二章 §2 第 2 目），都是平坦 A-模。
\item
  若 A 是半单环（《代数学》第八章 §5 第 1 目定义 1），则每个右 A-模 E 都是半单的，因而是一些单模的直和；由于这些单模中的每一个都同构于 \(\mathbf{A}_d\) 的一个直因子（出处同上，§5 第 1 目命题 6），E 是投射的，因此由 (1) 是平坦的（参看习题 16）。
\end{enumerate}

*(3) 在第二、三章中，我们将详细研究平坦 A-模的两个重要例子：分式环 \(S^{-1}A\) 以及 \(\hat{A}\) （即 \(A\) 关于 \(\mathfrak{I}\) -进拓扑的 Hausdorff 完备化）。*

命题 3. 设 A 是环，E 是右 A-模。

\begin{enumerate}
\def\labelenumi{(\roman{enumi})}
\item
  设 E 是平坦的。对 A 的每个不是 0 的右除子(*) 的元素 a，关系 \(x \in \mathbf{E}\) 、 \(xa = 0\) 蕴含 \(x = 0\) 。
\item
  设 A 是一个整环，其中每个有限生成理想都是主理想（例如主理想整环（《代数学》第七章 §1 第 1 目））。则 E 是平坦的，必要且充分的是 E 是无挠的。
\end{enumerate}

我们证明 (i)。设 \(v: \mathbf{A}_s \to \mathbf{A}_s\) 是左 A-模同态 \(t \mapsto ta\) ；假设表明 \(v\) 是单射的。由于 E 是平坦的，同态 \(1_{\mathbf{E}} \otimes v: \mathbf{E} \otimes_{\mathbf{A}} \mathbf{A}_s \to \mathbf{E} \otimes_{\mathbf{A}} \mathbf{A}_s\) 也是单射的。当 \(\mathbf{E} \otimes_{\mathbf{A}} \mathbf{A}_s\) 典范地等同于 E 时， \(1_{\mathbf{E}} \otimes v\) 就成为 E 的自同态 \(x \mapsto xa\) 。于是关系 \(xa = 0\) 蕴含 \(x = 0\) 。

我们证明 (ii)。由 (i)，若 E 是平坦的，则 E 是无挠的。反之，设 E 是无挠 A-模；我们验证，对 A 的每个有限生成理想 a，典范同态 \(E \otimes_{A} a \to E\) 是单射的（第 3 目，注记 1）。若 \(a = (0)\) ，这一断言是显然的；否则，由假设 a = Aa，其中 \(a \in A\) ， \(a \neq 0\) ，而 \(t \mapsto ta\) 就是 A 到 a 上的一个同构 v；用 i 表示典范单射 \(a \to A\) ， \(i \circ v\) 就是 E 上比为 a 的位似，且由于假设 E 无挠，它是单射的。于是 \(1_{E} \otimes (i \circ v) = (1_{E} \otimes i) \circ (1_{E} \otimes v)\) ；由于 \(1_{E} \otimes v\) 是同构， \(1_{E} \otimes i\) 是单射的，这就完成了证明。

例. 把命题 3 应用于环 \(\mathbf{Z}\) ，可以看出 \(\mathbf{Q}\) 是平坦 \(\mathbf{Z}\) -模，但 \(\mathbf{Z}/n\mathbf{Z}\) （对 \(n \geqslant 2\) ）不是平坦 \(\mathbf{Z}\) -模。

